% ===================================================================
%  图表第 7 号 — 冬天的村子。— 春天的农庄。
%  Traduction 2026 d'apres texto/io, controlee sur texto/fr, texto/en
%  et texto/es. Consigne : voir texto/zh/10-tubiao-01.tex.
% ===================================================================

%%K t07-noto-1 noto
\VUnotes{143.2mm}{%
(*) 一顿饭的细目见图表第 6 号和第 13 号。%
}

%%K t07-apar-1 apar
\VUsaut{4.0mm}
\VUcentre{13.2pt}{90}{第二组}
\VUsaut{3.0mm}
\VUfilet{16mm}
\VUsaut{5.6mm}
\VUtitre{15.0pt}{60}{图表第 7 号}
\VUsaut{3.0mm}

%%K t07-tit-1 sub
\VUcentre{12.0pt}{0}{\textit{第一场。}}
\VUsaut{3.0mm}

%%K t07-tit-2 sub
\VUpk{10.2pt}{40}{冬天的村子。}
\VUsaut{4.0mm}

%%K t07-c1-01-1 p
1. --- 新年假期里，我跟父亲去了新庄，那是个小村子，他在那里有事要办。

%%K t07-c1-02-1 p
2. --- 我们坐的是城里和那村子之间来往的\VUgras{驿车} \textsuperscript{(11)}；可是
天很冷，坐这么不舒服的车走一趟，一点儿也不好受。

%%K t07-c1-03-1 p
3. --- 所以当我们看见\VUgras{车夫}把车停在广场上、停在\textit{白熊}
\VUgras{客店} \textsuperscript{(2)} 门前时，真是高兴。\VUgras{店主}
\textsuperscript{(4)} 正在门口等我们，安安静静地抽着他的\VUgras{烟斗}。

%%K t07-c1-03-2 p
他请我们进去，我用不着人请第二遍，就在炉膛里噼啪响着的葡萄藤火前坐下了。
那时我就不在乎那刺骨的\VUgras{北风}了，那风吹得老\VUgras{招牌}
\textsuperscript{(3)} 吱吱作响，把烟囱里冒出来的\VUgras{烟} \textsuperscript{(1)} 卷成
一团团黑影带走。

%%K t07-c1-04-1 p
4. --- 正是吃饭的时候。我们不再耽搁，把端上来的那顿简单却可口的饭菜好好
吃了一顿(*)。

%%K t07-tit-3 sub
\VUpk{10.2pt}{40}{几次拜访。}
\VUsaut{3.0mm}

%%K t07-c1-05-1 p
5. --- 饭后我跟着做保险代理的父亲去拜访他的主顾。我们先去看住在客店隔壁
的\VUgras{铁匠} \textsuperscript{(5)}。他正在给一匹马钉掌。我很佩服他帮手的力气，
那人把马\VUgras{蹄}稳稳抵在自己的皮围裙上，铁匠就抡起大锤把\VUgras{马掌}
\textsuperscript{(6)} 钉上。

%%K t07-c1-06-1 p
6. --- 然后我们去村里的\VUgras{鞋匠} \textsuperscript{(7)} 老雅各那里。虽说他的
招牌是一只很气派的\VUgras{靴子}，他却只在给一双穿破的旧鞋换底。

%%K t07-c1-06-2 p
老头儿笑着对我们说：「在城里我是『鞋匠』，一个主顾也没有。在这儿我是
『补鞋的』，活多得做不完。」

%%K t07-c1-07-1 p
7. --- 他跟我们讲起他的邻居\VUgras{理发匠} \textsuperscript{(8)}，那人的主顾没有
一个满意的。他的\VUgras{剃刀}老是有豁口，\VUgras{剪子}从来剪不利索。可是
村里就他一个理发匠，人们自然只好让他刮脸剪头。他还是个又抠又难说话的人。

%%K t07-c1-08-1 p
8. --- 幸好别的\VUgras{邻居}不像他。譬如\VUgras{车匠} \textsuperscript{(9)} 就是个
好人，勤快又肯帮忙。把车交给他修是件痛快事。他的活儿总是收拾得干干净净。

%%K t07-c1-09-1 p
9. --- 老雅各对\VUgras{鞍匠} \textsuperscript{(10)} 也没有话说，活紧的时候他有时
还帮着缝\VUgras{马具}。

%%K t07-c1-09-2 p
父亲问起他家里人。「都好。我孙女在\VUgras{学校} \textsuperscript{(12)}。
\VUgras{女教师} \textsuperscript{(13)} 很喜欢她；她是个好\VUgras{学生}
\textsuperscript{(14)}。不像我那个淘气儿子；他心里只有玩。\VUgras{先生}
\textsuperscript{(15)} 拿他一点办法也没有，什么也教不进去。」

%%K t07-tit-4 sub
\VUsaut{3.0mm}
\VUpk{10.2pt}{40}{村里的闲话。}
\VUsaut{3.0mm}

%%K t07-c1-10-1 p
10. --- 接着老头儿讲起村里的新闻。要盖一所新学校了，因为\VUgras{村公所}里
那间早就太小了。这也不难，只要买下\VUgras{磨坊主}的一块园子就行。这对那
可怜人倒正合适。天一冷起来，\VUgras{磨坊} \textsuperscript{(16)} 的\VUgras{磨盘}
就磨不了麦子了，因为\VUgras{水轮} \textsuperscript{(17)} 被\VUgras{冰}
\textsuperscript{(25)} 冻住了，那是\VUgras{小溪} \textsuperscript{(23)} 的冰。

%%K t07-c1-11-1 p
11. --- 「说到盖房子，您看见我们新的\VUgras{教堂} \textsuperscript{(18)} 了吗？
盖着一层雪，看着很好看。\VUgras{乌鸦} \textsuperscript{(19)} 认不出自己的老钟楼
了，闷闷地落在\VUgras{墓地} \textsuperscript{(20)} 那棵大树上。」

%%K t07-c1-12-1 p
12. --- 说到墓地，鞋匠就想起父亲认识的、这一年里过世的人。「我的先生，
又添了多少座\VUgras{坟} \textsuperscript{(21)} 啊，就在那\VUgras{柏树}
\textsuperscript{(22)} 底下！我们的老公证人也走了。全村都去送\VUgras{葬}。」

%%K t07-c1-13-1 p
13. --- 「他的接班人正在溪上滑冰。他刚才还来过，叫我给他修\VUgras{冰鞋}
\textsuperscript{(26)} 断了的带子。」

%%K t07-c1-14-1 p
14. --- 「那边溪岸上是新来的\VUgras{村警} \textsuperscript{(27)}。他从前是宪兵，
是村里那些淘气鬼的克星。他们一看见他的\VUgras{绑腿} \textsuperscript{(29)} 和
\VUgras{制帽} \textsuperscript{(28)} 就跑。」

%%K t07-c1-15-1 p
15. --- 「啊！老约瑟给面包师送\VUgras{一车柴捆} \textsuperscript{(30)} 来了。---
不错，父亲说，我认得他那队\VUgras{骡子} \textsuperscript{(31)}。我正有话要跟他说，
这就去。再会了，老雅各。」

%%K t07-c1-16-1 p
16. --- 我们正要出门，村里的孩子放了学，一窝蜂涌到\VUgras{冰道}
\textsuperscript{(33)} 上，也不怕\VUgras{摔} \textsuperscript{(34)} 断胳膊腿。其中
一个从车匠那儿取来他的\VUgras{雪橇} \textsuperscript{(32)}，让一个朋友坐上去，
拉着就跑。

%%K t07-c1-17-1 p
17. --- 另一些孩子冲到一处\VUgras{雪堆} \textsuperscript{(37)} 那里，就在
\VUgras{桥} \textsuperscript{(40)} 边，动手打他们头一天堆的\VUgras{雪人} \textsuperscript{(39)}，
他们给雪人戴了帽子、插了烟斗，还给它挎上了书包。

%%K t07-c1-18-1 p
18. --- 他们不在乎冰冷的风和严寒。多半人连条\VUgras{围巾} \textsuperscript{(35)}
都没有；打完\VUgras{雪仗} \textsuperscript{(38)} 要暖和暖和，只要向老雅克琳
\VUgras{小贩}买一把烫手的\VUgras{栗子} \textsuperscript{(42)} 就够了；那老太太
裹着一条太薄的\VUgras{披巾} \textsuperscript{(43)}，冻得直哆嗦。

%%K t07-tit-5 sub
\VUsaut{3.0mm}
\VUcentre{12.0pt}{0}{\textit{第二场。}}
\VUsaut{3.0mm}

%%K t07-tit-6 sub
\VUpk{10.2pt}{40}{春天的农庄。}
\VUsaut{4.0mm}

%%K t07-c2-01-1 p
1. --- 冬天固然有它的乐趣，我们还是满心欢喜地迎接\VUgras{春天}的归来。

%%K t07-c2-01-2 p
五月里傍晚的农庄院子，是多么叫人高兴的一幅画啊！

%%K t07-c2-02-1 p
2. --- 天边\VUgras{太阳} \textsuperscript{(1)} 慢慢落下去，\VUgras{田野}
\textsuperscript{(3)} 沐在它绯红的\VUgras{光} \textsuperscript{(2)} 里。勤快的\VUgras{犁地人}
\textsuperscript{(6)} 正赶着犁他的\VUgras{地} \textsuperscript{(4)}。

%%K t07-c2-02-2 p
他的\VUgras{犁} \textsuperscript{(5)} 由一匹壮\VUgras{马} \textsuperscript{(7)}
拉着，划出\VUgras{垄沟} \textsuperscript{(8)}，\VUgras{播种人} \textsuperscript{(9)} 就要一把把
往里撒下\VUgras{种子}，将来结成金黄的穗子。

%%K t07-c2-03-1 p
3. --- \VUgras{割草人} \textsuperscript{(11)} 用镰刀割了草场的草，翻晒的人把
\VUgras{干草} \textsuperscript{(10)} 用\VUgras{草叉} \textsuperscript{(12)} 和耙子
翻晒干了，这会儿正往家走。

%%K t07-c2-03-2 p
有的走在牛拉的重车前头，肩上扛着家伙。有的懒些，舒舒服服地躺在香喷喷的
干草上。

%%K t07-c2-04-1 p
4. --- 走过时他们亲热地跟\VUgras{园丁} \textsuperscript{(16)} 打招呼，那园丁正在
\VUgras{果园} \textsuperscript{(13)} 里忙着剪\VUgras{果树}、给\VUgras{梨树}
\textsuperscript{(14)} 嫁接。\VUgras{李树} \textsuperscript{(15)} 开着花，上足了肥的
\VUgras{菜园} \textsuperscript{(17)} 里，菜蔬、白菜和莴苣也已经长得很好。

%%K t07-c2-04-2 p
矮墙上一只漂亮的\VUgras{孔雀} \textsuperscript{(18)} 开着屏，显摆它华丽的羽毛。

%%K t07-tit-7 sub
\VUpk{10.2pt}{40}{农庄的院子。}
\VUsaut{3.0mm}

%%K t07-c2-05-1 p
5. --- \VUgras{马厩} \textsuperscript{(19)} 的门开着，看得见草料架和
\VUgras{垫草} \textsuperscript{(23)}，那是\VUgras{母马} \textsuperscript{(21)} 的，
\VUgras{马夫} \textsuperscript{(20)} 刚给它卸了套。\VUgras{小马驹} \textsuperscript{(22)} 长腿细细的
样子很滑稽，一蹦一跳地跟着母亲。

%%K t07-c2-06-1 p
6. --- \VUgras{井} \textsuperscript{(29)} 边\VUgras{母牛} \textsuperscript{(25)} 走过来
喝水，喝的是那只当作水槽的木\VUgras{槽} \textsuperscript{(26)}。\VUgras{小牛}
\textsuperscript{(24)} 趁机吃奶，\VUgras{公牛} \textsuperscript{(27)} 闻到草料的香味，
快活地哞哞叫着，骄傲地扬起长着有力\VUgras{角} \textsuperscript{(28)} 的头。

%%K t07-c2-07-1 p
7. --- 人人都在干活。\VUgras{女工} \textsuperscript{(32)} 拉着井\VUgras{绳}
\textsuperscript{(31)}。她用\VUgras{水桶} \textsuperscript{(30)} 打水给牲口喝。
\VUgras{谷仓} \textsuperscript{(33)} 旁边有个人在耙散落的草，\VUgras{农舍}
\textsuperscript{(34)} 门前一个老\VUgras{纺妇} \textsuperscript{(35)} 守着纺车和
\VUgras{线杆}，像从前一样纺\VUgras{亚麻}或\VUgras{大麻}。

%%K t07-tit-8 sub
\VUsaut{3.0mm}
\VUpk{10.2pt}{40}{农庄主。}
\VUsaut{3.0mm}

%%K t07-c2-08-1 p
8. --- \VUgras{农庄主} \textsuperscript{(36)} 指挥着这一切活计，给伙计们派活。他
吩咐把\VUgras{耙} \textsuperscript{(40)} 和\VUgras{十字镐} \textsuperscript{(39)} 收进
棚里去，那两样丢在\VUgras{窝} \textsuperscript{(38)} 旁边了，就是那条\VUgras{狗}
\textsuperscript{(37)} 的窝。

%%K t07-c2-09-1 p
9. --- 他的吆喝惊起一只\VUgras{鸽子} \textsuperscript{(41)}，它拼命飞回
\VUgras{鸽舍} \textsuperscript{(42)}；也惊了\VUgras{母鸡} \textsuperscript{(43)}，它们
急忙躲回\VUgras{鸡窝} \textsuperscript{(44)}。

%%K t07-c2-10-1 p
10. --- \VUgras{羊圈} \textsuperscript{(48)} 前面，老\VUgras{牧人} \textsuperscript{(45)}
正赶回他的\VUgras{羊群} \textsuperscript{(49)}。他拿着\VUgras{牧杖}
\textsuperscript{(46)}，那杖跟他褪了色的\VUgras{罩衫} \textsuperscript{(47)} 一样从不
离身；他把\VUgras{公绵羊} \textsuperscript{(50)}、\VUgras{母绵羊} \textsuperscript{(53)}、
\VUgras{羔羊} \textsuperscript{(54)}、\VUgras{山羊} \textsuperscript{(51)} 和
\VUgras{小山羊} \textsuperscript{(52)} 都赶进圈里。他忠心的\VUgras{狗}
\textsuperscript{(55)} 阿耳戈斯走在最后，汪汪叫着催那些掉队的。

%%K t07-c2-11-1 p
11. --- 这一片吵闹忙乱并不搅扰那头好\VUgras{奶牛} \textsuperscript{(56)} 的安静，
它一边摇响自己的\VUgras{铃铛} \textsuperscript{(57)}，一边在\VUgras{牛栏}
\textsuperscript{(58)} 门前让人挤奶。

%%K t07-c2-12-1 p
12. --- 它像是懂得自己该出奶，好让\VUgras{农妇} \textsuperscript{(60)} 在
\VUgras{搅乳器} \textsuperscript{(61)} 里打\VUgras{黄油}，或者做那些上好的
\VUgras{奶酪} \textsuperscript{(64)}——奶酪正搁在横在\VUgras{木桶} \textsuperscript{(63)}
上的板子上沥水。

%%K t07-c2-13-1 p
13. --- 整个\VUgras{奶房} \textsuperscript{(59)} 都由\VUgras{长工} \textsuperscript{(65)}
收拾得很干净，他刚在手里提的大桶里洗了\VUgras{奶罐} \textsuperscript{(62)}。

%%K t07-tit-9 sub
\VUsaut{3.0mm}
\VUpk{10.2pt}{40}{家禽院。}
\VUsaut{3.0mm}

%%K t07-c2-14-1 p
14. --- 他穿着笨重的木鞋走得挺蹒跚，于是惊跑了\VUgras{火鸡} \textsuperscript{(68)}、
\VUgras{母鸭} \textsuperscript{(66)}、\VUgras{公鸭} \textsuperscript{(67)} 和\VUgras{鹅}
\textsuperscript{(69)}，它们张大\VUgras{嘴} \textsuperscript{(70)} 惊惶地叫成一片。

%%K t07-c2-14-2 p
\VUgras{公鸡} \textsuperscript{(71)} 更好斗些，昂起红\VUgras{冠} \textsuperscript{(72)}，
抖了抖\VUgras{尾巴} \textsuperscript{(73)} 上的毛，放开嗓子「喔喔喔」地叫了一声。

%%K t07-c2-15-1 p
15. --- 一只\VUgras{老母鸡} \textsuperscript{(74)} 在\VUgras{粪堆}
\textsuperscript{(81)} 上刨食，带着它的\VUgras{小鸡} \textsuperscript{(75)}。\VUgras{小猪}
\textsuperscript{(78)} 朝母亲跑去，那是猪圈旁看得见的一头大\VUgras{母猪}
\textsuperscript{(77)}。

%%K t07-c2-16-1 p
16. --- 笼子里的\VUgras{兔子} \textsuperscript{(79)} 正起劲地吃农庄主的女儿给它们
带来的嫩\VUgras{草} \textsuperscript{(80)}。让妮很疼她的兔子，每天从鸡窝里取了
新鲜的\VUgras{蛋} \textsuperscript{(76)} 以后，总要来看看她这些小朋友。
\VUsaut{6.0mm}
\VUfilet{22mm}
